%------------------------------------------------------------------------------%
\begin{question}
  Determine os autovalores da matriz
  \(
    A =
    \begin{bmatrix}
      4 & 1 \\
      2 & 3
    \end{bmatrix}
  \)
\end{question}

%------------------------------------------------------------------------------%
\begin{resolution}{Polinômio Característica}
  \begin{align*}
    \onslide<+->{p(\lambda)}
    \onslide<+->{& = \det(A - \lambda I)}
    \\
    \onslide<+->{& = \det\left(
          \begin{bmatrix}
            4 & 1 \\
            2 & 3
          \end{bmatrix}
          - \lambda
          \begin{bmatrix}
            1 & 0 \\
            0 & 1
          \end{bmatrix}
    \right)}
    \\
    \onslide<+->{& = \begin{vmatrix}[c]
          4-\lambda & 1         \\
          2         & 3-\lambda
        \end{vmatrix}}
    \\
    \onslide<+->{& = (4-\lambda)(3-\lambda) - 1 \times 2  \\}
    \onslide<+->{& = 12 - 4\lambda - 3\lambda + \lambda^2 - 2        \\}
    \onslide<+->{& = \lambda^2 - 7\lambda + 10}
  \end{align*}
\end{resolution}

%------------------------------------------------------------------------------%
\begin{resolution}{Equação Característica}
  \[
    p(\lambda)
    \;=\; \lambda^2 - 7\lambda + 10
    \pause
    \;=\; 0
  \]

  \[
    \pause
    \Delta
    \pause
    \;=\; b^2 - 4ac
    \pause
    \;=\; (-7)^2 - 4 \times 1 \times 10
    \pause
    \;=\; 49 - 40
    \pause
    \;=\; 9
  \]

  \[
    \pause
    \sqrt{\Delta}
    \pause
    \;=\; \sqrt{9}
    \pause
    \;=\; 3
  \]

  \[
    \pause
    \lambda_1
    \pause
    \;=\; \frac{-b + \sqrt{\Delta}}{2a}
    \pause
    \;=\; \frac{7 + 3}{2}
    \pause
    \;=\; 5
  \]

  \[
    \pause
    \lambda_2
    \pause
    \;=\; \frac{-b - \sqrt{\Delta}}{2a}
    \pause
    \;=\; \frac{7 - 3}{2}
    \pause
    \;=\; 2
  \]
\end{resolution}

%------------------------------------------------------------------------------%
\endinput
